\documentclass{syllabus}

\begin{document}

\thispagestyle{fancy}
\coursecode{PHY 2970}
\date{Spring 2026}
\author{}
\title{\coursecode: {Portfolio Seminar A}}
\maketitle

\meeting{ Tuesday: & 8:00 AM - 9:40 AM, DSC 054 }
\officehours{TR: & 11:30 AM - 12:30 PM}
\prerequisites{PHY 2210:& (General Physics II, co-enrolled)}


\section*{General Information}

\infotableport





\section*{Course overview}

\textbf{Catalog description:} The first course in the Physics Portfolio Seminar concentrates on building student experience and expertise in using physics principles to solve problems and approach complex scenarios. Emphasis is on group work and the informal and formal presentations of solutions to open-ended questions. Students will also participate in a handful of career and professional development exercises and be introduced to the concept of a physics portfolio. Prior completion or concurrent enrollment in PHY 2210 is required.

The purpose of this course is to get hands-on experience with solving more open-ended physics problems. Points of emphasis will include:
\begin{enumerate}
    \item Identifying principles and strategies to solve complex problems
    \item Use of basic computer simulation to answer complex questions,
    \item Design of simple experiments using sensor data to answer complex physics questions.
\end{enumerate}

One theme of the problem-solving elements of the course will be exploring deviations from the idealized behavior presented in the general physics sequence. Students will create original work and present it in written and oral formats.

\subsection*{Student learning outcomes} 
\noindent

Upon successful completion of this course, students will be able to:
\begin{enumerate}[label=\textcolor{CarthageDark}{\bfseries\arabic*.}]
    \item Solve challenging problems using fundamental principles that span more than one subfield of physics.
    \item Create quantitative models of real-world phenomena and apply them as part of the problem-solving process.
    \item Create computer programs and/or use spreadsheet software to perform the quantitative analyses that arise in problem solving.
    \item Students will use appropriate physics vocabulary and concepts in written and oral reports on outcomes of problem solving.
\end{enumerate}



\section*{Course components}

\paragraph{Challenge problems:} In the first section of the course, students will work on weekly challenge problems, which need to be formulated into quantitative models which are then deployed, often as basic computer calculations. 

Each week, students will informally explain their work on the challenge problem from the previous week to the instructor.

\paragraph{Experiment design:} In the second section of the course, students will design experiments using available sensors to determine some form of deviation from idealized behavior. This will involve data collection and data analysis.

Students will give verbal updates each week on their progress on the experiment.

\paragraph{Portfolio elements:} Students will create a more formal written report and oral presentation on one of their challenge problems and one of their experiments in the final weeks of the course.

\section*{Grading Scheme}
\subsection*{Component Weights}

\noindent\begin{tabular}{lr}
\textbf{Component} & \textbf{Weight}       \\ \hline
\textbf{Challenges}      & 35\%  \\
\textbf{Experiments}           & 35\%       \\
\textbf{Portfolio}            & 30\%     
\end{tabular}%



\lettergrades{a) demonstrated effort in class, b) willingness to seek help outside of class.}


\tipsforsuccess{problem}

\section*{Policies}
\subsection{General course policies}
Expect a response within one business day for emails sent to the instructor. To ensure a prompt response, format the subject line as: [\coursecode] *insert subject here*

All required material will be covered in lecture, but can also be found in the parallel readings, so in the case of absence students should consult the scheduled readings. Students found sleeping will be brought to the attention of the class for public ridicule and no effort made to wake them up.

Students are expected to arrive and leave on time. I will attempt to arrive five minutes early and be available for at least five minutes after class.

\latework{None}

\collegepolicies[Other college policies]

\end{document}
